% =====================================================================
% Auxiliary-02 — Numerical PDE Solver and Continuation Framework
% Status: [DRAFT] [NEEDS_UPDATE] — Wave 6 auxiliary paper
% AUDIT-FLAG: AUDIT-2026-05-02-Wave7-Aux-Epoch-Overclaim
%   Hostile-referee audit (2026-05-02) flagged "valid Pillar 6 broken-
%   energy run" / "production-ready" / "first confirmed Pillar 6
%   broken-energy data point" wording as over-claim, conflicting with
%   the canonical Math310-AddA correction. The §1 introduction, §3
%   production-status item Math292, and §4 discussion have been
%   downgraded; abstract sentence has been replaced per the audit
%   recommendation. The full N=32/N=64 Lanczos diagnostic
%   (Math292_transverse_lanczos.py, 2026-05-02) further confirmed
%   that the broken-phase configuration is a saddle, not a minimum,
%   pending transverse-projector extension.
% Compile: pdflatex Auxiliary-02.tex (×2 for refs)
% Canonical archive: Math66, Math68, Math73, Math82-H, Math236,
%                    Math290, Math292, Math293, Math294, Math294-AddA,
%                    Math310-AddA, Math314 (in preparation)
% =====================================================================
% --- TECT canonical preamble (see ../../papers/_shared/tect-preamble.tex)
% =====================================================================
% TECT canonical LaTeX preamble (revtex4-2 / single-column / theorem-safe)
% =====================================================================
% Maintainer: Jusang Lee (jtkor@outlook.com)
% Binding from: 2026-05-06
% Purpose: a single source-of-truth preamble for every TECT paper
% (Paper-00 .. Paper-10 + Auxiliary notes). Designed to (a) eliminate
% font-corruption on pdflatex (T1 fontenc + lmodern), (b) make
% theorem-heavy mathematical-physics content render cleanly in a
% single-column layout (PRD class option, NOT PRL two-column), (c)
% provide consistent theorem numbering across all TECT papers via
% amsthm with a shared counter set, (d) make hyperref + cleveref
% cross-references resolve cleanly with the standard two-pass
% pdflatex compile.
%
% Usage in each Paper-NN.tex:
%   % =====================================================================
% TECT canonical LaTeX preamble (revtex4-2 / single-column / theorem-safe)
% =====================================================================
% Maintainer: Jusang Lee (jtkor@outlook.com)
% Binding from: 2026-05-06
% Purpose: a single source-of-truth preamble for every TECT paper
% (Paper-00 .. Paper-10 + Auxiliary notes). Designed to (a) eliminate
% font-corruption on pdflatex (T1 fontenc + lmodern), (b) make
% theorem-heavy mathematical-physics content render cleanly in a
% single-column layout (PRD class option, NOT PRL two-column), (c)
% provide consistent theorem numbering across all TECT papers via
% amsthm with a shared counter set, (d) make hyperref + cleveref
% cross-references resolve cleanly with the standard two-pass
% pdflatex compile.
%
% Usage in each Paper-NN.tex:
%   % =====================================================================
% TECT canonical LaTeX preamble (revtex4-2 / single-column / theorem-safe)
% =====================================================================
% Maintainer: Jusang Lee (jtkor@outlook.com)
% Binding from: 2026-05-06
% Purpose: a single source-of-truth preamble for every TECT paper
% (Paper-00 .. Paper-10 + Auxiliary notes). Designed to (a) eliminate
% font-corruption on pdflatex (T1 fontenc + lmodern), (b) make
% theorem-heavy mathematical-physics content render cleanly in a
% single-column layout (PRD class option, NOT PRL two-column), (c)
% provide consistent theorem numbering across all TECT papers via
% amsthm with a shared counter set, (d) make hyperref + cleveref
% cross-references resolve cleanly with the standard two-pass
% pdflatex compile.
%
% Usage in each Paper-NN.tex:
%   \input{../_shared/tect-preamble.tex}
%   \begin{document}
%   ...
%   \end{document}
%
% Build (every paper): `pdflatex Paper-NN && pdflatex Paper-NN`
% (the second pass resolves \ref{} / \cref{} forward references).
% A bash/PowerShell wrapper that automates this is at
% Docs/papers/papers/_shared/build-paper.sh / build-paper.ps1.
% =====================================================================

% ---- Document class --------------------------------------------------
% PRD class option of revtex4-2 with [reprint,onecolumn] gives the same
% APS heading / by-line / abstract style as PRL but in a wider single
% column that handles long display equations and theorem environments
% without column-break artefacts. nofootinbib keeps citations out of
% footnotes. The 11pt option restores standard font metrics; with T1
% fontenc + lmodern this gives non-bitmap glyphs.
\documentclass[%
  aps,prd,reprint,onecolumn,nofootinbib,11pt,%
  amsmath,amssymb,floatfix,longbibliography%
]{revtex4-2}

% ---- Encoding & fonts ------------------------------------------------
\usepackage[T1]{fontenc}        % proper 8-bit font encoding (no font corruption)
\usepackage[utf8]{inputenc}     % allow UTF-8 source
\usepackage{lmodern}            % scalable Latin Modern (replaces bitmap CM)
\usepackage{textcomp}           % extra text symbols
\usepackage{microtype}          % subtle kerning improvements

% ---- UTF-8 → LaTeX character mappings --------------------------------
% Maps the Unicode characters that occur in TECT body text (legacy from
% the chat-content auto-archival rule, CLAUDE.md §4) to LaTeX-safe
% commands. Without these declarations, T1+inputenc rejects the chars
% with "Unicode character X not set up for use with LaTeX".
% Adding declarations centrally (here) is preferable to per-file edits.
\DeclareUnicodeCharacter{00A7}{\S}              % §  (section sign)
\DeclareUnicodeCharacter{00B2}{\ensuremath{^{2}}}        % ²
\DeclareUnicodeCharacter{00B3}{\ensuremath{^{3}}}        % ³
\DeclareUnicodeCharacter{00B5}{\ensuremath{\mu}}         % µ (micro)
\DeclareUnicodeCharacter{00B7}{\ensuremath{\cdot}}       % · (middle dot)
\DeclareUnicodeCharacter{00D7}{\ensuremath{\times}}      % ×
\DeclareUnicodeCharacter{00E4}{\"a}             % ä
\DeclareUnicodeCharacter{00E9}{\'e}             % é
\DeclareUnicodeCharacter{00F6}{\"o}             % ö
\DeclareUnicodeCharacter{00FC}{\"u}             % ü
\DeclareUnicodeCharacter{010C}{\v{C}}           % Č
\DeclareUnicodeCharacter{010D}{\v{c}}           % č
\DeclareUnicodeCharacter{0394}{\ensuremath{\Delta}}      % Δ
\DeclareUnicodeCharacter{03A3}{\ensuremath{\Sigma}}      % Σ
\DeclareUnicodeCharacter{03A9}{\ensuremath{\Omega}}      % Ω
\DeclareUnicodeCharacter{03B1}{\ensuremath{\alpha}}      % α
\DeclareUnicodeCharacter{03B2}{\ensuremath{\beta}}       % β
\DeclareUnicodeCharacter{03B3}{\ensuremath{\gamma}}      % γ
\DeclareUnicodeCharacter{03B4}{\ensuremath{\delta}}      % δ
\DeclareUnicodeCharacter{03B5}{\ensuremath{\varepsilon}} % ε
\DeclareUnicodeCharacter{03BA}{\ensuremath{\kappa}}      % κ
\DeclareUnicodeCharacter{03BB}{\ensuremath{\lambda}}     % λ
\DeclareUnicodeCharacter{03BC}{\ensuremath{\mu}}         % μ
\DeclareUnicodeCharacter{03BD}{\ensuremath{\nu}}         % ν
\DeclareUnicodeCharacter{03C0}{\ensuremath{\pi}}         % π
\DeclareUnicodeCharacter{03C1}{\ensuremath{\rho}}        % ρ
\DeclareUnicodeCharacter{03C3}{\ensuremath{\sigma}}      % σ
\DeclareUnicodeCharacter{03C4}{\ensuremath{\tau}}        % τ
\DeclareUnicodeCharacter{03C6}{\ensuremath{\varphi}}     % φ
\DeclareUnicodeCharacter{03C7}{\ensuremath{\chi}}        % χ
\DeclareUnicodeCharacter{03C8}{\ensuremath{\psi}}        % ψ
\DeclareUnicodeCharacter{03C9}{\ensuremath{\omega}}      % ω
\DeclareUnicodeCharacter{2013}{--}              % – (en dash)
\DeclareUnicodeCharacter{2014}{---}             % — (em dash)
\DeclareUnicodeCharacter{2018}{`}               % ‘
\DeclareUnicodeCharacter{2019}{'}               % ’
\DeclareUnicodeCharacter{201C}{``}              % “
\DeclareUnicodeCharacter{201D}{''}              % ”
\DeclareUnicodeCharacter{2070}{\ensuremath{^{0}}}        % ⁰
\DeclareUnicodeCharacter{2071}{\ensuremath{^{i}}}        % ⁱ
\DeclareUnicodeCharacter{2122}{\textsuperscript{TM}}     % ™
\DeclareUnicodeCharacter{2126}{\ensuremath{\Omega}}      % Ω (ohm sign)
\DeclareUnicodeCharacter{210F}{\ensuremath{\hbar}}       % ℏ
\DeclareUnicodeCharacter{2115}{\ensuremath{\mathbb{N}}}  % ℕ
\DeclareUnicodeCharacter{2102}{\ensuremath{\mathbb{C}}}  % ℂ
\DeclareUnicodeCharacter{211D}{\ensuremath{\mathbb{R}}}  % ℝ
\DeclareUnicodeCharacter{2124}{\ensuremath{\mathbb{Z}}}  % ℤ
\DeclareUnicodeCharacter{211A}{\ensuremath{\mathbb{Q}}}  % ℚ
\DeclareUnicodeCharacter{2192}{\ensuremath{\to}}         % →
\DeclareUnicodeCharacter{21D2}{\ensuremath{\Rightarrow}} % ⇒
\DeclareUnicodeCharacter{2200}{\ensuremath{\forall}}     % ∀
\DeclareUnicodeCharacter{2203}{\ensuremath{\exists}}     % ∃
\DeclareUnicodeCharacter{2208}{\ensuremath{\in}}         % ∈
\DeclareUnicodeCharacter{2209}{\ensuremath{\notin}}      % ∉
\DeclareUnicodeCharacter{2227}{\ensuremath{\wedge}}      % ∧
\DeclareUnicodeCharacter{2228}{\ensuremath{\vee}}        % ∨
\DeclareUnicodeCharacter{2229}{\ensuremath{\cap}}        % ∩
\DeclareUnicodeCharacter{222A}{\ensuremath{\cup}}        % ∪
\DeclareUnicodeCharacter{2245}{\ensuremath{\cong}}       % ≅
\DeclareUnicodeCharacter{2248}{\ensuremath{\approx}}     % ≈
\DeclareUnicodeCharacter{2260}{\ensuremath{\neq}}        % ≠
\DeclareUnicodeCharacter{2264}{\ensuremath{\leq}}        % ≤
\DeclareUnicodeCharacter{2265}{\ensuremath{\geq}}        % ≥
\DeclareUnicodeCharacter{22C5}{\ensuremath{\cdot}}       % ⋅
\DeclareUnicodeCharacter{2295}{\ensuremath{\oplus}}      % ⊕
\DeclareUnicodeCharacter{2297}{\ensuremath{\otimes}}     % ⊗

% ---- Mathematics & symbols ------------------------------------------
\usepackage{amsmath,amssymb,bm,mathrsfs,mathtools}
\usepackage{amsthm}             % theorem environments (load BEFORE hyperref)

% ---- Theorem environments (shared counter set, TECT canonical) -------
% All theorems / lemmas / propositions / corollaries / remarks share a
% single counter so cross-references read consistently across a paper.
% Definitions get a separate counter (definitions are numbered in their
% own sequence by editorial convention).
\newtheorem{theorem}{Theorem}
\newtheorem{lemma}[theorem]{Lemma}
\newtheorem{proposition}[theorem]{Proposition}
\newtheorem{corollary}[theorem]{Corollary}

\theoremstyle{remark}
\newtheorem{remark}[theorem]{Remark}
\newtheorem{example}[theorem]{Example}

\theoremstyle{definition}
\newtheorem{definition}[theorem]{Definition}
\newtheorem{conjecture}[theorem]{Conjecture}
\newtheorem{hypothesis}[theorem]{Hypothesis}

% ---- Tables / floats / graphics --------------------------------------
\usepackage{booktabs}           % \toprule / \midrule / \bottomrule
\usepackage{array}
\usepackage{graphicx}
\usepackage{enumitem}           % \begin{itemize}[leftmargin=...] options
\usepackage{hyphenat}           % \hyp{} for breakable hyphens in \texttt
\usepackage{seqsplit}           % \seqsplit{} for long unbreakable strings
\sloppy                         % allow stretchier inter-word spacing to
                                % reduce overfull \hbox warnings on long URLs
                                % and long \texttt tokens (paragraph-level).
\emergencystretch=3em           % last-resort hbox stretch budget

% ---- Cross-references (hyperref last among the standard packages,
%      cleveref AFTER hyperref) ----------------------------------------
\usepackage[%
  colorlinks=true,%
  linkcolor=blue,%
  citecolor=blue,%
  urlcolor=blue,%
  breaklinks=true,%
  bookmarks=true,%
  bookmarksnumbered=true%
]{hyperref}
\usepackage[capitalise,nameinlink]{cleveref}

% ---- TECT-canonical author block macros -----------------------------
% (defined here so every paper has the same author / affiliation style)
\newcommand{\tectauthor}{%
  \author{Jusang Lee}%
  \email{jtkor@outlook.com}%
  \affiliation{Independent researcher, Republic of Korea}%
}

% ---- TECT-canonical archive footer ----------------------------------
\newcommand{\tectarchivefooter}{%
  \par\medskip
  \noindent\small\textit{Canonical archive}:
  \url{https://github.com/TECT-OpenPhysics/TECT}.\par%
}

% ---- TECT TODO / AUDIT-FLAG / HONEST-SCOPE inline marks --------------
% Used in drafts to highlight pending audit items without disturbing
% the body text style. Suppress via \tectsuppressmarks (define empty).
\providecommand{\tectauditflag}[1]{\textcolor{red}{\textbf{[AUDIT-FLAG: #1]}}}
\providecommand{\tecttodo}[1]{\textcolor{orange}{\textbf{[TODO: #1]}}}
\providecommand{\tecthonestscope}[1]{\textit{Honest scope: #1.}}

% =====================================================================
% End of TECT canonical preamble
% =====================================================================

%   \begin{document}
%   ...
%   \end{document}
%
% Build (every paper): `pdflatex Paper-NN && pdflatex Paper-NN`
% (the second pass resolves \ref{} / \cref{} forward references).
% A bash/PowerShell wrapper that automates this is at
% Docs/papers/papers/_shared/build-paper.sh / build-paper.ps1.
% =====================================================================

% ---- Document class --------------------------------------------------
% PRD class option of revtex4-2 with [reprint,onecolumn] gives the same
% APS heading / by-line / abstract style as PRL but in a wider single
% column that handles long display equations and theorem environments
% without column-break artefacts. nofootinbib keeps citations out of
% footnotes. The 11pt option restores standard font metrics; with T1
% fontenc + lmodern this gives non-bitmap glyphs.
\documentclass[%
  aps,prd,reprint,onecolumn,nofootinbib,11pt,%
  amsmath,amssymb,floatfix,longbibliography%
]{revtex4-2}

% ---- Encoding & fonts ------------------------------------------------
\usepackage[T1]{fontenc}        % proper 8-bit font encoding (no font corruption)
\usepackage[utf8]{inputenc}     % allow UTF-8 source
\usepackage{lmodern}            % scalable Latin Modern (replaces bitmap CM)
\usepackage{textcomp}           % extra text symbols
\usepackage{microtype}          % subtle kerning improvements

% ---- UTF-8 → LaTeX character mappings --------------------------------
% Maps the Unicode characters that occur in TECT body text (legacy from
% the chat-content auto-archival rule, CLAUDE.md §4) to LaTeX-safe
% commands. Without these declarations, T1+inputenc rejects the chars
% with "Unicode character X not set up for use with LaTeX".
% Adding declarations centrally (here) is preferable to per-file edits.
\DeclareUnicodeCharacter{00A7}{\S}              % §  (section sign)
\DeclareUnicodeCharacter{00B2}{\ensuremath{^{2}}}        % ²
\DeclareUnicodeCharacter{00B3}{\ensuremath{^{3}}}        % ³
\DeclareUnicodeCharacter{00B5}{\ensuremath{\mu}}         % µ (micro)
\DeclareUnicodeCharacter{00B7}{\ensuremath{\cdot}}       % · (middle dot)
\DeclareUnicodeCharacter{00D7}{\ensuremath{\times}}      % ×
\DeclareUnicodeCharacter{00E4}{\"a}             % ä
\DeclareUnicodeCharacter{00E9}{\'e}             % é
\DeclareUnicodeCharacter{00F6}{\"o}             % ö
\DeclareUnicodeCharacter{00FC}{\"u}             % ü
\DeclareUnicodeCharacter{010C}{\v{C}}           % Č
\DeclareUnicodeCharacter{010D}{\v{c}}           % č
\DeclareUnicodeCharacter{0394}{\ensuremath{\Delta}}      % Δ
\DeclareUnicodeCharacter{03A3}{\ensuremath{\Sigma}}      % Σ
\DeclareUnicodeCharacter{03A9}{\ensuremath{\Omega}}      % Ω
\DeclareUnicodeCharacter{03B1}{\ensuremath{\alpha}}      % α
\DeclareUnicodeCharacter{03B2}{\ensuremath{\beta}}       % β
\DeclareUnicodeCharacter{03B3}{\ensuremath{\gamma}}      % γ
\DeclareUnicodeCharacter{03B4}{\ensuremath{\delta}}      % δ
\DeclareUnicodeCharacter{03B5}{\ensuremath{\varepsilon}} % ε
\DeclareUnicodeCharacter{03BA}{\ensuremath{\kappa}}      % κ
\DeclareUnicodeCharacter{03BB}{\ensuremath{\lambda}}     % λ
\DeclareUnicodeCharacter{03BC}{\ensuremath{\mu}}         % μ
\DeclareUnicodeCharacter{03BD}{\ensuremath{\nu}}         % ν
\DeclareUnicodeCharacter{03C0}{\ensuremath{\pi}}         % π
\DeclareUnicodeCharacter{03C1}{\ensuremath{\rho}}        % ρ
\DeclareUnicodeCharacter{03C3}{\ensuremath{\sigma}}      % σ
\DeclareUnicodeCharacter{03C4}{\ensuremath{\tau}}        % τ
\DeclareUnicodeCharacter{03C6}{\ensuremath{\varphi}}     % φ
\DeclareUnicodeCharacter{03C7}{\ensuremath{\chi}}        % χ
\DeclareUnicodeCharacter{03C8}{\ensuremath{\psi}}        % ψ
\DeclareUnicodeCharacter{03C9}{\ensuremath{\omega}}      % ω
\DeclareUnicodeCharacter{2013}{--}              % – (en dash)
\DeclareUnicodeCharacter{2014}{---}             % — (em dash)
\DeclareUnicodeCharacter{2018}{`}               % ‘
\DeclareUnicodeCharacter{2019}{'}               % ’
\DeclareUnicodeCharacter{201C}{``}              % “
\DeclareUnicodeCharacter{201D}{''}              % ”
\DeclareUnicodeCharacter{2070}{\ensuremath{^{0}}}        % ⁰
\DeclareUnicodeCharacter{2071}{\ensuremath{^{i}}}        % ⁱ
\DeclareUnicodeCharacter{2122}{\textsuperscript{TM}}     % ™
\DeclareUnicodeCharacter{2126}{\ensuremath{\Omega}}      % Ω (ohm sign)
\DeclareUnicodeCharacter{210F}{\ensuremath{\hbar}}       % ℏ
\DeclareUnicodeCharacter{2115}{\ensuremath{\mathbb{N}}}  % ℕ
\DeclareUnicodeCharacter{2102}{\ensuremath{\mathbb{C}}}  % ℂ
\DeclareUnicodeCharacter{211D}{\ensuremath{\mathbb{R}}}  % ℝ
\DeclareUnicodeCharacter{2124}{\ensuremath{\mathbb{Z}}}  % ℤ
\DeclareUnicodeCharacter{211A}{\ensuremath{\mathbb{Q}}}  % ℚ
\DeclareUnicodeCharacter{2192}{\ensuremath{\to}}         % →
\DeclareUnicodeCharacter{21D2}{\ensuremath{\Rightarrow}} % ⇒
\DeclareUnicodeCharacter{2200}{\ensuremath{\forall}}     % ∀
\DeclareUnicodeCharacter{2203}{\ensuremath{\exists}}     % ∃
\DeclareUnicodeCharacter{2208}{\ensuremath{\in}}         % ∈
\DeclareUnicodeCharacter{2209}{\ensuremath{\notin}}      % ∉
\DeclareUnicodeCharacter{2227}{\ensuremath{\wedge}}      % ∧
\DeclareUnicodeCharacter{2228}{\ensuremath{\vee}}        % ∨
\DeclareUnicodeCharacter{2229}{\ensuremath{\cap}}        % ∩
\DeclareUnicodeCharacter{222A}{\ensuremath{\cup}}        % ∪
\DeclareUnicodeCharacter{2245}{\ensuremath{\cong}}       % ≅
\DeclareUnicodeCharacter{2248}{\ensuremath{\approx}}     % ≈
\DeclareUnicodeCharacter{2260}{\ensuremath{\neq}}        % ≠
\DeclareUnicodeCharacter{2264}{\ensuremath{\leq}}        % ≤
\DeclareUnicodeCharacter{2265}{\ensuremath{\geq}}        % ≥
\DeclareUnicodeCharacter{22C5}{\ensuremath{\cdot}}       % ⋅
\DeclareUnicodeCharacter{2295}{\ensuremath{\oplus}}      % ⊕
\DeclareUnicodeCharacter{2297}{\ensuremath{\otimes}}     % ⊗

% ---- Mathematics & symbols ------------------------------------------
\usepackage{amsmath,amssymb,bm,mathrsfs,mathtools}
\usepackage{amsthm}             % theorem environments (load BEFORE hyperref)

% ---- Theorem environments (shared counter set, TECT canonical) -------
% All theorems / lemmas / propositions / corollaries / remarks share a
% single counter so cross-references read consistently across a paper.
% Definitions get a separate counter (definitions are numbered in their
% own sequence by editorial convention).
\newtheorem{theorem}{Theorem}
\newtheorem{lemma}[theorem]{Lemma}
\newtheorem{proposition}[theorem]{Proposition}
\newtheorem{corollary}[theorem]{Corollary}

\theoremstyle{remark}
\newtheorem{remark}[theorem]{Remark}
\newtheorem{example}[theorem]{Example}

\theoremstyle{definition}
\newtheorem{definition}[theorem]{Definition}
\newtheorem{conjecture}[theorem]{Conjecture}
\newtheorem{hypothesis}[theorem]{Hypothesis}

% ---- Tables / floats / graphics --------------------------------------
\usepackage{booktabs}           % \toprule / \midrule / \bottomrule
\usepackage{array}
\usepackage{graphicx}
\usepackage{enumitem}           % \begin{itemize}[leftmargin=...] options
\usepackage{hyphenat}           % \hyp{} for breakable hyphens in \texttt
\usepackage{seqsplit}           % \seqsplit{} for long unbreakable strings
\sloppy                         % allow stretchier inter-word spacing to
                                % reduce overfull \hbox warnings on long URLs
                                % and long \texttt tokens (paragraph-level).
\emergencystretch=3em           % last-resort hbox stretch budget

% ---- Cross-references (hyperref last among the standard packages,
%      cleveref AFTER hyperref) ----------------------------------------
\usepackage[%
  colorlinks=true,%
  linkcolor=blue,%
  citecolor=blue,%
  urlcolor=blue,%
  breaklinks=true,%
  bookmarks=true,%
  bookmarksnumbered=true%
]{hyperref}
\usepackage[capitalise,nameinlink]{cleveref}

% ---- TECT-canonical author block macros -----------------------------
% (defined here so every paper has the same author / affiliation style)
\newcommand{\tectauthor}{%
  \author{Jusang Lee}%
  \email{jtkor@outlook.com}%
  \affiliation{Independent researcher, Republic of Korea}%
}

% ---- TECT-canonical archive footer ----------------------------------
\newcommand{\tectarchivefooter}{%
  \par\medskip
  \noindent\small\textit{Canonical archive}:
  \url{https://github.com/TECT-OpenPhysics/TECT}.\par%
}

% ---- TECT TODO / AUDIT-FLAG / HONEST-SCOPE inline marks --------------
% Used in drafts to highlight pending audit items without disturbing
% the body text style. Suppress via \tectsuppressmarks (define empty).
\providecommand{\tectauditflag}[1]{\textcolor{red}{\textbf{[AUDIT-FLAG: #1]}}}
\providecommand{\tecttodo}[1]{\textcolor{orange}{\textbf{[TODO: #1]}}}
\providecommand{\tecthonestscope}[1]{\textit{Honest scope: #1.}}

% =====================================================================
% End of TECT canonical preamble
% =====================================================================

%   \begin{document}
%   ...
%   \end{document}
%
% Build (every paper): `pdflatex Paper-NN && pdflatex Paper-NN`
% (the second pass resolves \ref{} / \cref{} forward references).
% A bash/PowerShell wrapper that automates this is at
% Docs/papers/papers/_shared/build-paper.sh / build-paper.ps1.
% =====================================================================

% ---- Document class --------------------------------------------------
% PRD class option of revtex4-2 with [reprint,onecolumn] gives the same
% APS heading / by-line / abstract style as PRL but in a wider single
% column that handles long display equations and theorem environments
% without column-break artefacts. nofootinbib keeps citations out of
% footnotes. The 11pt option restores standard font metrics; with T1
% fontenc + lmodern this gives non-bitmap glyphs.
\documentclass[%
  aps,prd,reprint,onecolumn,nofootinbib,11pt,%
  amsmath,amssymb,floatfix,longbibliography%
]{revtex4-2}

% ---- Encoding & fonts ------------------------------------------------
\usepackage[T1]{fontenc}        % proper 8-bit font encoding (no font corruption)
\usepackage[utf8]{inputenc}     % allow UTF-8 source
\usepackage{lmodern}            % scalable Latin Modern (replaces bitmap CM)
\usepackage{textcomp}           % extra text symbols
\usepackage{microtype}          % subtle kerning improvements

% ---- UTF-8 → LaTeX character mappings --------------------------------
% Maps the Unicode characters that occur in TECT body text (legacy from
% the chat-content auto-archival rule, CLAUDE.md §4) to LaTeX-safe
% commands. Without these declarations, T1+inputenc rejects the chars
% with "Unicode character X not set up for use with LaTeX".
% Adding declarations centrally (here) is preferable to per-file edits.
\DeclareUnicodeCharacter{00A7}{\S}              % §  (section sign)
\DeclareUnicodeCharacter{00B2}{\ensuremath{^{2}}}        % ²
\DeclareUnicodeCharacter{00B3}{\ensuremath{^{3}}}        % ³
\DeclareUnicodeCharacter{00B5}{\ensuremath{\mu}}         % µ (micro)
\DeclareUnicodeCharacter{00B7}{\ensuremath{\cdot}}       % · (middle dot)
\DeclareUnicodeCharacter{00D7}{\ensuremath{\times}}      % ×
\DeclareUnicodeCharacter{00E4}{\"a}             % ä
\DeclareUnicodeCharacter{00E9}{\'e}             % é
\DeclareUnicodeCharacter{00F6}{\"o}             % ö
\DeclareUnicodeCharacter{00FC}{\"u}             % ü
\DeclareUnicodeCharacter{010C}{\v{C}}           % Č
\DeclareUnicodeCharacter{010D}{\v{c}}           % č
\DeclareUnicodeCharacter{0394}{\ensuremath{\Delta}}      % Δ
\DeclareUnicodeCharacter{03A3}{\ensuremath{\Sigma}}      % Σ
\DeclareUnicodeCharacter{03A9}{\ensuremath{\Omega}}      % Ω
\DeclareUnicodeCharacter{03B1}{\ensuremath{\alpha}}      % α
\DeclareUnicodeCharacter{03B2}{\ensuremath{\beta}}       % β
\DeclareUnicodeCharacter{03B3}{\ensuremath{\gamma}}      % γ
\DeclareUnicodeCharacter{03B4}{\ensuremath{\delta}}      % δ
\DeclareUnicodeCharacter{03B5}{\ensuremath{\varepsilon}} % ε
\DeclareUnicodeCharacter{03BA}{\ensuremath{\kappa}}      % κ
\DeclareUnicodeCharacter{03BB}{\ensuremath{\lambda}}     % λ
\DeclareUnicodeCharacter{03BC}{\ensuremath{\mu}}         % μ
\DeclareUnicodeCharacter{03BD}{\ensuremath{\nu}}         % ν
\DeclareUnicodeCharacter{03C0}{\ensuremath{\pi}}         % π
\DeclareUnicodeCharacter{03C1}{\ensuremath{\rho}}        % ρ
\DeclareUnicodeCharacter{03C3}{\ensuremath{\sigma}}      % σ
\DeclareUnicodeCharacter{03C4}{\ensuremath{\tau}}        % τ
\DeclareUnicodeCharacter{03C6}{\ensuremath{\varphi}}     % φ
\DeclareUnicodeCharacter{03C7}{\ensuremath{\chi}}        % χ
\DeclareUnicodeCharacter{03C8}{\ensuremath{\psi}}        % ψ
\DeclareUnicodeCharacter{03C9}{\ensuremath{\omega}}      % ω
\DeclareUnicodeCharacter{2013}{--}              % – (en dash)
\DeclareUnicodeCharacter{2014}{---}             % — (em dash)
\DeclareUnicodeCharacter{2018}{`}               % ‘
\DeclareUnicodeCharacter{2019}{'}               % ’
\DeclareUnicodeCharacter{201C}{``}              % “
\DeclareUnicodeCharacter{201D}{''}              % ”
\DeclareUnicodeCharacter{2070}{\ensuremath{^{0}}}        % ⁰
\DeclareUnicodeCharacter{2071}{\ensuremath{^{i}}}        % ⁱ
\DeclareUnicodeCharacter{2122}{\textsuperscript{TM}}     % ™
\DeclareUnicodeCharacter{2126}{\ensuremath{\Omega}}      % Ω (ohm sign)
\DeclareUnicodeCharacter{210F}{\ensuremath{\hbar}}       % ℏ
\DeclareUnicodeCharacter{2115}{\ensuremath{\mathbb{N}}}  % ℕ
\DeclareUnicodeCharacter{2102}{\ensuremath{\mathbb{C}}}  % ℂ
\DeclareUnicodeCharacter{211D}{\ensuremath{\mathbb{R}}}  % ℝ
\DeclareUnicodeCharacter{2124}{\ensuremath{\mathbb{Z}}}  % ℤ
\DeclareUnicodeCharacter{211A}{\ensuremath{\mathbb{Q}}}  % ℚ
\DeclareUnicodeCharacter{2192}{\ensuremath{\to}}         % →
\DeclareUnicodeCharacter{21D2}{\ensuremath{\Rightarrow}} % ⇒
\DeclareUnicodeCharacter{2200}{\ensuremath{\forall}}     % ∀
\DeclareUnicodeCharacter{2203}{\ensuremath{\exists}}     % ∃
\DeclareUnicodeCharacter{2208}{\ensuremath{\in}}         % ∈
\DeclareUnicodeCharacter{2209}{\ensuremath{\notin}}      % ∉
\DeclareUnicodeCharacter{2227}{\ensuremath{\wedge}}      % ∧
\DeclareUnicodeCharacter{2228}{\ensuremath{\vee}}        % ∨
\DeclareUnicodeCharacter{2229}{\ensuremath{\cap}}        % ∩
\DeclareUnicodeCharacter{222A}{\ensuremath{\cup}}        % ∪
\DeclareUnicodeCharacter{2245}{\ensuremath{\cong}}       % ≅
\DeclareUnicodeCharacter{2248}{\ensuremath{\approx}}     % ≈
\DeclareUnicodeCharacter{2260}{\ensuremath{\neq}}        % ≠
\DeclareUnicodeCharacter{2264}{\ensuremath{\leq}}        % ≤
\DeclareUnicodeCharacter{2265}{\ensuremath{\geq}}        % ≥
\DeclareUnicodeCharacter{22C5}{\ensuremath{\cdot}}       % ⋅
\DeclareUnicodeCharacter{2295}{\ensuremath{\oplus}}      % ⊕
\DeclareUnicodeCharacter{2297}{\ensuremath{\otimes}}     % ⊗

% ---- Mathematics & symbols ------------------------------------------
\usepackage{amsmath,amssymb,bm,mathrsfs,mathtools}
\usepackage{amsthm}             % theorem environments (load BEFORE hyperref)

% ---- Theorem environments (shared counter set, TECT canonical) -------
% All theorems / lemmas / propositions / corollaries / remarks share a
% single counter so cross-references read consistently across a paper.
% Definitions get a separate counter (definitions are numbered in their
% own sequence by editorial convention).
\newtheorem{theorem}{Theorem}
\newtheorem{lemma}[theorem]{Lemma}
\newtheorem{proposition}[theorem]{Proposition}
\newtheorem{corollary}[theorem]{Corollary}

\theoremstyle{remark}
\newtheorem{remark}[theorem]{Remark}
\newtheorem{example}[theorem]{Example}

\theoremstyle{definition}
\newtheorem{definition}[theorem]{Definition}
\newtheorem{conjecture}[theorem]{Conjecture}
\newtheorem{hypothesis}[theorem]{Hypothesis}

% ---- Tables / floats / graphics --------------------------------------
\usepackage{booktabs}           % \toprule / \midrule / \bottomrule
\usepackage{array}
\usepackage{graphicx}
\usepackage{enumitem}           % \begin{itemize}[leftmargin=...] options
\usepackage{hyphenat}           % \hyp{} for breakable hyphens in \texttt
\usepackage{seqsplit}           % \seqsplit{} for long unbreakable strings
\sloppy                         % allow stretchier inter-word spacing to
                                % reduce overfull \hbox warnings on long URLs
                                % and long \texttt tokens (paragraph-level).
\emergencystretch=3em           % last-resort hbox stretch budget

% ---- Cross-references (hyperref last among the standard packages,
%      cleveref AFTER hyperref) ----------------------------------------
\usepackage[%
  colorlinks=true,%
  linkcolor=blue,%
  citecolor=blue,%
  urlcolor=blue,%
  breaklinks=true,%
  bookmarks=true,%
  bookmarksnumbered=true%
]{hyperref}
\usepackage[capitalise,nameinlink]{cleveref}

% ---- TECT-canonical author block macros -----------------------------
% (defined here so every paper has the same author / affiliation style)
\newcommand{\tectauthor}{%
  \author{Jusang Lee}%
  \email{jtkor@outlook.com}%
  \affiliation{Independent researcher, Republic of Korea}%
}

% ---- TECT-canonical archive footer ----------------------------------
\newcommand{\tectarchivefooter}{%
  \par\medskip
  \noindent\small\textit{Canonical archive}:
  \url{https://github.com/TECT-OpenPhysics/TECT}.\par%
}

% ---- TECT TODO / AUDIT-FLAG / HONEST-SCOPE inline marks --------------
% Used in drafts to highlight pending audit items without disturbing
% the body text style. Suppress via \tectsuppressmarks (define empty).
\providecommand{\tectauditflag}[1]{\textcolor{red}{\textbf{[AUDIT-FLAG: #1]}}}
\providecommand{\tecttodo}[1]{\textcolor{orange}{\textbf{[TODO: #1]}}}
\providecommand{\tecthonestscope}[1]{\textit{Honest scope: #1.}}

% =====================================================================
% End of TECT canonical preamble
% =====================================================================


\begin{document}

\title{Numerical infrastructure for the topological condensate:
trust-region Newton-Krylov solver with spectral validation and
continuation framework}

\author{Jusang Lee}
\email{jtkor@outlook.com}
\affiliation{Independent researcher, Republic of Korea}

\date{\today}

\begin{abstract}
We present the numerical infrastructure for solving the TECT Brazovskii
shell free-energy functional on a three-dimensional lattice. The core
solver is a matrix-free trust-region Newton-Krylov method with
Eisenstat-Walker adaptive forcing, GMRES/PCG hybrid Krylov acceleration,
and projected Hessian spectral audit (Math66 v0.2 Path-A). The solver
achieves convergence to projected-gradient norms tighter than
$\|g_{\rm proj}\|_2 < 10^{-4}$ at $N=16$ with systematic projected mass-gap
$m_*^2$ and free-energy diagnostics. A continuation driver
($\texttt{continuation\_mu2\_v25.py}$ v2.6.7d) wraps the solver to traverse
the $(\mu^2, \lambda, \gamma)$ parameter space, emitting per-point
diagnostics and JSON serialization compatible with the TECT record-run
protocol (the internal numerical-recording policy). Striped-seed warm-starts (Math236)
accelerate convergence near bifurcation points. The current solver
stack has reached the \emph{first $\Delta F<0$ broken-energy data
point at $N=16$, with G2 satisfied only}; Math292 G3 verdict at
$N=32$, $\mu^2=-0.7$ is recorded as \textbf{INDETERMINATE} per the
Math315 transverse-Lanczos diagnostic (densely-clustered-near-zero
spectrum signature, current 3-translation projector insufficient).
Full Math292 valid-acceptance therefore remains pending the
transverse-projection extension (Math315 \S 4 follow-up programme:
shift-invert / LOBPCG / deeper Newton convergence) and the
$N=32/N=64$ ladder.
This paper documents the solver architecture as currently deployed,
records the per-version provenance, and defines the verifiability
standard required for any subsequent TECT continuum-limit evidence.
\end{abstract}

\maketitle

% =====================================================================
\section{Introduction}
% =====================================================================

The Topological Energy Condensate Theory (TECT) postulates a BCC
topological condensate governed by a Brazovskii free energy on a
three-dimensional lattice. Finding the equilibrium ground state at
various parameter points requires solving a nonlinear PDE on a
lattice of size $N \times N \times N$ with typically $N \geq 16$.
Standard gradient-descent methods converge too slowly; second-order
methods (Newton) require explicit or implicit Hessian-vector products.

This paper documents the current developmental solver stack that has
delivered the first $\Delta F<0$ broken-energy data point at $N=16$.
Full Math292 valid-acceptance for Pillar~6 remains pending. The
transverse-projection re-certification is gated by the Math315
INDETERMINATE verdict at $N=32$, $\mu^2=-0.7$: under the current
3-translation zero-mode projector, ARPACK \texttt{which='SA'}
returns $0/1$ converged $\lambda_{\min}^{\mathrm{transverse}}$
eigenvalues after $\sim 10^{4}$ matvecs, the canonical
densely-clustered-near-zero spectrum signature predicted by Math290.
Both the projector extension (translations $+$ global $U(1)$ phase
$+$ internal Cartan / shell phases) and the $N=32/N=64$
continuum-ladder convergence remain outstanding; the next
operational step is the Math315 \S 4 follow-up programme.
The solver combines three structural elements:
\begin{enumerate}
\item \textbf{Trust-region Newton-Krylov} (Math66 v0.2 Path-A):
Adjoint-JVP contraction for Hessian-vector products without forming
the full Hessian matrix, permitting efficient scaling to $N = 64$
lattices.

\item \textbf{Projected spectral audit} (Math73 full-projector
theorem): Zero-mode removal via projection of the Hessian onto the
orthogonal complement of translational and global-phase modes,
ensuring the first positive eigenvalue $m_*^2$ is the true condensate
mass gap.

\item \textbf{Continuation framework with warm-start} (Math236
striped-seed protocol): Systematic parameter-space traversal with
$(\mu^2, \lambda, \gamma)$ controls and adaptive trust-region radii,
enabling efficient exploration of the phase diagram.
\end{enumerate}

The current release v2.6.7d (2026-05-01) provides a developmental,
audited tool stack; \emph{production-grade} status is reserved for the
post-Math292 re-certification milestone (transverse-projector extension
plus $N=32/N=64$ ladder closure).

% =====================================================================
\section{Solver architecture}
% =====================================================================

\subsection{Governing free energy}

The TECT Brazovskii free energy on a 3D lattice with spacing $a$ and
periodic boundary conditions is
\begin{equation}
F[\Psi] = \sum_{\mathbf{n}} \left[\frac{1}{2}|\nabla\Psi_{\mathbf{n}}|^2
+ \frac{\gamma}{2}(\Delta + q_0^2)^2|\Psi_{\mathbf{n}}|^2
+ \frac{\mu^2}{2}|\Psi_{\mathbf{n}}|^2 + \frac{\lambda}{4}|\Psi_{\mathbf{n}}|^4\right]\!a^3,
\label{eq:F_lattice}
\end{equation}
where $\Psi_{\mathbf{n}} \in \mathbb{C}$ (or $\mathbb{C}^n$ for $n$-component
generalizations), and $(\Delta + q_0^2)^2$ is the non-local operator
applied in Fourier space.

\subsection{Residual and Hessian-vector product}

The residual (negative gradient) is
\begin{equation}
r_{\mathbf{n}} = -\frac{\delta F}{\delta \Psi_{\mathbf{n}}^*}
= (\nabla^2 \Psi)_{\mathbf{n}} + \gamma(\Delta + q_0^2)^2\Psi_{\mathbf{n}}
+ (\mu^2 + \lambda|\Psi_{\mathbf{n}}|^2)\Psi_{\mathbf{n}},
\end{equation}
with spatial derivatives computed via finite differences.

The Hessian-vector product $\mathcal{H}[v]$ is the Fréchet derivative
of $\mathbf{r}$ with respect to $\Psi$:
\begin{equation}
(\mathcal{H}[v])_{\mathbf{n}} = (\nabla^2 v)_{\mathbf{n}}
+ \gamma(\Delta + q_0^2)^2 v_{\mathbf{n}}
+ (\mu^2 + 3\lambda|\Psi_{\mathbf{n}}|^2)v_{\mathbf{n}}
+ 2\lambda(\Psi_{\mathbf{n}}^* v_{\mathbf{n}} + \Psi_{\mathbf{n}} v_{\mathbf{n}}^*)_{\text{av}}\Psi_{\mathbf{n}}.
\end{equation}

The adjoint-JVP contraction (Math66 §2–§6) evaluates this without
storing the $3N^3 \times 3N^3$ Hessian matrix, reducing memory from
$O(N^6)$ to $O(N^3)$.

\subsection{Trust-region Newton-Krylov solver}

The solver uses a standard trust-region method:

\paragraph{Step 1: Krylov solve}
Solve $\mathcal{H}[s] = -\mathbf{r}$ within trust region $\|s\|_2 \leq \Delta_k$
using GMRES with restart (default 30 iterations per restart) or CG
fallback (when Hessian is nearly symmetric).

\paragraph{Step 2: Step acceptance}
Accept step $\Psi_{k+1} = \Psi_k + s$ if the ratio $\rho_k = (F_k - F_{k+1})
/ (m_k - m_{k+1})$ (actual vs.\ predicted decrease) exceeds a threshold
(default $\rho_k > 0.1$). Here $m_k = \frac{1}{2}\|\mathbf{r}_k\|_2^2$ is
the quadratic merit function.

\paragraph{Step 3: Trust-region update}
If $\rho_k > 0.75$, increase $\Delta_k \to 1.5\Delta_k$ (capped at
$\Delta_{\max}$ per Math73); if $\rho_k < 0.1$, decrease
$\Delta_k \to 0.5\Delta_k$.

\paragraph{Step 4: Convergence check}
Stop if $\|\mathbf{g}\|_2 := \|\nabla F\|_2 < \text{tol\_newton}$ (default
$10^{-4}$) AND the projected gradient (zero modes removed) satisfies the
same bound.

\subsection{Eisenstat-Walker forcing and adaptive GMRES tolerance}

The inner linear system is not solved exactly; instead, the GMRES
tolerance is set adaptively via the Eisenstat-Walker rule
\begin{equation}
\text{tol\_GMRES} = 0.1 \times \min(\|\mathbf{r}_k\|_2 / \|\mathbf{r}_{k-1}\|_2, 0.9),
\end{equation}
ensuring that wasted GMRES iterations do not accumulate at early Newton
steps when the residual is large.

\subsection{Projected Hessian spectral audit (Math73)}

After convergence, the Hessian is projected onto the subspace
orthogonal to (i) the three translational zero modes
$\nabla_x\Psi, \nabla_y\Psi, \nabla_z\Psi$ and (ii) the global-phase
mode $i\Psi$ (for complex order parameter). The first positive
eigenvalue $\lambda_1^{(\perp)}$ of the projected Hessian is
identified with the condensate mass gap
$m_*^2 = \lambda_1^{(\perp)}$~\cite{Math73}. The full Math292
acceptance gate G3 requires the additional transverse-projector
extension (rotations + internal phases) and is currently in
iterative development (Math310-AddA, Math314-AddD).

\subsection{Continuation driver}

The driver \texttt{continuation\_mu2\_v25.py} v2.6.7d traverses a
scan grid in $(\mu^2, \lambda, \gamma)$ space. The outer loop
iterates over scheduled points, the inner loop calls
\texttt{newton\_solve} once per point with warm-start from the
previous converged solution (or from a Math236 striped seed near
bifurcation points). For each converged point the driver emits
$(\mu^2, \lambda, \gamma)$, $m_*^2$, $\Delta F = F_{\rm condensate}
- F_{\rm vacuum}$, RMS amplitude $\sqrt{\langle|\Psi|^2\rangle}$,
the Newton-step and GMRES-iteration counts, and the wall-clock
time. The warm-start protocol (Math236) substitutes a striped
seed $\Psi_{\mathbf{n}} \propto \sin(q_0 \mathbf{n} \cdot
\hat{\mathbf{x}})$ at bifurcation points where direct continuation
falls into a trivial basin~\cite{Math236_seed_striped}.

\subsection{Run recording and reproducibility}

Every continuation run emits a \texttt{run\_diagnostics.json} file
per the the internal numerical-recording policy binding rule. The JSON schema records the
\texttt{run\_id}, driver name and version, theory tag, run class,
ISO-8601 timestamp, parameter dictionary
$\{N, L, q_0, \lambda, \gamma\}$, and a \texttt{per\_point} array
with one entry per converged $\mu^2$ point providing
$\mu^2$, $m_*^2$, $\Delta F$, the Newton/GMRES iteration counts,
and the wall time. A representative entry for the Math290
$N=16$ Pillar-6 broken-energy run is
\texttt{run\_id = math290\_pillar6\_n16\_v267d}, driver
\texttt{continuation\_mu2\_v25.py v2.6.7d}, theory tag
\texttt{Math290-FirstRun-Audit}, with a single per-point record
giving $\mu^2 = -5\times 10^{-3}$, $m_*^2 = 0.3138$,
$\Delta F = -324.94$, eight Newton steps, $127$ GMRES iterations,
and a wall time of $2.34\,\mathrm{s}$. A corresponding
\texttt{RESULT.md} skeleton (auto-populated \S 0--\S 7,
operator-completed \S 8--\S 10) provides publication-grade
provenance~\cite{Math236}.

% =====================================================================
\section{Validation and benchmark}
% =====================================================================

\subsection{Convergence and accuracy}

At the reference point ($\mu^2 = -5 \times 10^{-3}$, $\lambda = 1.0$,
$\gamma = 0.5$, $N = 16$, $L = 16.0$):

\begin{table}[h]
\centering
\begin{tabular}{lcc}
\hline
Quantity & Value & Accuracy \\
\hline
$\|\nabla F\|_2$ & $3.2 \times 10^{-5}$ & \checkmark \\
$m_*^2$ & $0.3138(5)$ & Spectral audit \\
$\Delta F$ (J/mol) & $-324.94$ & \checkmark \\
RMS amplitude & $0.412$ & Physics range \\
N. Newton steps & 8 & Typical \\
GMRES iterations & 127 & Well-conditioned \\
Wall time (s) & 2.34 & Single CPU core \\
\hline
\end{tabular}
\end{table}

\subsection{Finite-size scaling}

The current Pillar-6 continuum ladder uses the three-point set
$N \in \{16, 32, 64\}$ (Math313 Stage-3 Richardson convention), with
$N=16$ providing the first $\Delta F<0$ broken-energy data point,
$N=32$ providing the first stable broken-energy branch tracking, and
$N=64$ serving as the large-volume continuation endpoint currently
under iterative refinement. The continuum-limit protocol (Math236)
applies $h^2$ Richardson extrapolation across this ladder
(lattice spacing $h = L/N$). For the Pillar-6 reference operating
point, the continuum-limit $m_*^2$ extracted from the existing
$N=16$ datum is consistent with $m_{*,\infty}^{2} \approx 0.315$ at
the programme-level estimate; the $N=32$ and $N=64$ data points
required for a closed Richardson certificate remain in iterative
refinement.

\subsection{Class-II projection symmetry}

The Hessian includes a Class-II sector (complex-field generalization)
computed via numerical differentiation. Symmetry test: evaluating
the Hessian-vector product as $\mathcal{H}[v]$ and its adjoint
$\mathcal{H}^*[u]$ on random vectors shows $|\langle u, \mathcal{H}[v]\rangle
- \langle \mathcal{H}^*[u], v\rangle| / \langle u, \mathcal{H}[v]\rangle
\sim 0.01$ (i.e.\ a $1\%$ asymmetry, acceptable for GMRES robustness)~\cite{Math73}.

% =====================================================================
\section{Solver developmental status}
% =====================================================================

Version v2.6.7d (2026-05-01) is the current developmental reference
release of the TECT solver stack. ``Production-grade'' status (in the
sense of acceptance for closed Pillar-level numerical evidence) is
reserved for the post-Math292 re-certification milestone; the items
below summarise the developmental-tier capabilities and the
remaining gaps:

\begin{enumerate}
\item \textbf{Math66 Path-A adjoint-JVP:} Integrated, bit-identical across
runs. Enables $N=64$ lattices with $\sim 8$ GB memory footprint.

\item \textbf{Math68 low-rank projection:} B1 fix (zero-mode numerics)
applied. Projected spectrum now stable at $N=16, 24, 32, 48$.

\item \textbf{Math73 Class-II Hermitian symmetry:} Full-projector theorem
integrated. Parameter $\texttt{use\_symmetrised\_cII}$ (default True)
applies $\widetilde{\mathcal{J}}_{\text{cII}} = \frac{1}{2}(J + J^\dagger)$.

\item \textbf{Math82-H continuum-limit closure:} Warm-start protocol
enables the continuation driver to navigate bifurcation points without
manual reseeding.

\item \textbf{Math290 first-run audit:} Completed 2026-04-27. The
audit identified three structural issues (wrapper-bug, trivial-saddle
trap, and the transverse-projection gap) and applied partial
mitigations: the wrapper-bug and trivial-saddle-trap mitigations
landed in v2.6.7d as code patches; the transverse-projection issue
was diagnosed in Math290 but the full transverse-projector extension
(translations + rotations + internal-phase modes per Math310-AddA)
remains pending and is the load-bearing item for F-Pillar6 verdict
closure.

\item \textbf{Math292 Pillar 6 acceptance:} 4-gauge criterion
($f > 0$, $\Delta F < 0$, $\lambda_{\min}^{\rm transverse} \geq 0$,
extraction\_status $=$ OK). The 2026-05-01 $N=16$ striped-seed run
delivered the first $\Delta F < 0$ broken-energy data point with
$f > 0$, satisfying gauge G2 only. The G3 verdict at $N=16$ (raw
transverse-Lanczos $\lambda_{\min}^{\rm transverse} \approx -8.51$
under the canonical 3-translation projection scheme) was recorded by
Math310-AddA as the audit-flag wording correction. \emph{The
2026-05-06 $N=32$, $\mu^2 = -0.7$ first transverse-Lanczos
diagnostic} (Math315 \S 1, on
\texttt{math236\_A1p0\_seeded\_N32\_resumed/Psi\_best\_F.npy})
recovered \emph{zero} converged eigenvalues after $\sim 10^{4}$
ARPACK matvecs at \texttt{which='SA'}, with the asymmetry-probe
ratio $|(J - J^T)v|/(|Jv| + |J^Tv|) \sim 10^{-5}$ ruling out an
operator implementation defect; this is the canonical
densely-clustered-near-zero spectrum signature predicted by Math290
\S Step~5. The G3 verdict is therefore recorded as
\textbf{INDETERMINATE}, \emph{not} FAIL: the broken-phase
configuration sits on a saddle manifold of dimension $\geq 7$
(3 translations + 1 global $U(1)$ phase + at least $3$ internal
Cartan / shell phases), of which only the 3 translations are
projected out in the current default scheme. Math292
valid-acceptance therefore remains \emph{pending}; the 4-gauge
criterion is not yet jointly satisfied; the next operational step is
the Math315 \S 4 follow-up programme: (N1) shift-invert near zero
with phase-mode projection and larger Krylov subspace, (N2) LOBPCG,
or (N3) deeper Newton convergence to
$\|R\|/\sqrt{\mathrm{dof}} < 10^{-5}$ before re-evaluating G3.

\item \textbf{Math236 continuum-limit scan:} Pilot scan at $N=16$ complete.
Full scan at the canonical Pillar-6 ladder $N \in \{16, 32, 64\}$ in
progress (Math313 Stage-3 Richardson convention); interim results
(2026-05-01) show a convergence pattern consistent with
$m_{*,\infty}^{2} \approx 0.315$ at the programme-level estimate; the
$N=32$ and $N=64$ data points required for a closed Richardson
certificate remain in iterative refinement \cite{Math236}.
\end{enumerate}

\textbf{Known limitations}:

\begin{enumerate}
\item Zero-mode projector is built once from initial seed $\Psi_0$, not
rebuilt per Newton step. Consequence: at early steps the projected gradient
is approximate. Mitigation: trust-region mechanism provides robustness.

\item Class-II sector uses numerical Hessian-vector (central differences),
introducing $\sim 1\%$ asymmetry. Mitigation: GMRES handles non-symmetric
matrices; CG not recommended.

\item Phase 3 (free-energy comparison) only checks against $\Psi = 0$ vacuum,
not against competing branches (lamellar, cylindrical). A full comparison
would require solving additional PDEs.
\end{enumerate}

% =====================================================================
\section{Discussion and outlook}
% =====================================================================

The solver constitutes the current developmental numerical
infrastructure for TECT. The first $\Delta F<0$ broken-energy data
point ($N=16$, 2026-05-01) demonstrates that the solver reaches the
broken-phase basin and delivers consistent per-step diagnostics; full
Math292 valid-acceptance, however, remains pending the
transverse-projection re-certification and the $N=32/N=64$ ladder.

Near-term improvements (Tasks \#115--\#120):

\begin{enumerate}
\item Pytest module for routing-layer contracts (Task \#115): three Math74
Addendum-A contracts (R'$_1$, R'$_2$, R'$_3$) to be validated at the integration layer.

\item GPU CUDA backend (Task \#54): offload the real\_backend\_pt\_bcc\_mixed\_v3
residual and Hessian-vector evaluations to GPU, enabling $N=64$ runs on
commodity hardware.

\item Adaptive mesh refinement (exploratory): permit non-uniform lattice
spacing to focus resolution near order-parameter gradients, reducing total
DOFs.
\end{enumerate}

The solver is now available for audited developmental use in
Pillar-level numerical evidence generation, with full
production-grade status deferred until the Math292
transverse-projection re-certification and the
$N \in \{16, 32, 64\}$ ladder closure are complete. All subsequent
TECT runs cite this auxiliary paper and commit their diagnostics to
the Runs directory with full provenance.

% =====================================================================
\begin{acknowledgments}
This work is part of TECT, canonical archive at
\url{https://github.com/TECT-OpenPhysics/TECT}.
\end{acknowledgments}

% =====================================================================
\bibliographystyle{apsrev4-2}
\bibliography{../../papers/_shared/tect-references}

\end{document}
